% Lösungen Einheit 1 – Terme aufstellen und berechnen (zusammenbau v1.2)
\kbabschnitt{Terme aufstellen und berechnen}

\lbloes{21.}{a)}{$17$}
\lbloes{}{b)}{$3$}
\lbloes{}{c)}{$16$}
\lbloes{22.}{a)}{$4x - 3$}
\lbloes{}{b)}{$4x - 5$}
\lbloes{}{c)}{$3x + 8$}
\lbloes{}{d)}{$5x$}
\lbloes{}{e)}{$2x + 18$}
\lbloes{}{f)}{$2 \cdot (3x + 2)$}
\lbloes{23.}{a)}{$2 \cdot (3x - 7)$}
\lbloes{}{b)}{$6x - 2$}
\lbloes{24.}{a)}{Ein Kinobesuch kostet 15 € Eintritt, dazu kommen $x$ Getränke zu je 3 €; $x$ ist die Zahl der Getränke}
\lbloes{}{b)}{Lara hat 20 € und kauft $x$ Hefte zu je 2 €; $x$ ist die Zahl der Hefte}
\lbloes{}{c)}{Ein Fahrradverleih nimmt 10 € Grundgebühr und 4 € je Stunde; $x$ ist die Zahl der Stunden}
\lbloes{25.}{a)}{Zeile 2 ist falsch: Verdoppelt wird die ganze Summe, die Klammer fehlt.}
\lbloes{}{b)}{(1) falsch, z. B. $x = 10$: $10 - 3 = 7$, aber $3 - 10 = -7$. (2) wahr, denn ohne Klammer würde nur das erste Glied verdoppelt. (3) falsch, denn $2 \cdot (x + 4) = 2x + 8$ ist für jede Zahl $x$ um 4 größer als $2x + 4$.}
\lbloes{}{c)}{$4 \cdot (x + 6)$}
\lbloes{}{d)}{$56$ €}
